% !TEX root = ../../../main.tex
% Block card: interrogation.event_families -- "Event families" (Part I, the blocks chapter).
\cardfield{What it does}
Cuts the same-length stretch of signal around every event and sorts the events into families
of similar shape, using the same sorting methods as the Library.

\cardfield{When to use it}
When you suspect a recording holds more than one kind of event and want to know whether each
kind keeps its own rhythm. Put \blockref{interrogation.event_shape} before it and
\blockref{interrogation.intervals} after it: with \emph{write\_group} on (the default),
Intervals then times each family separately. The block is marked experimental: on real
recordings its families beat scrambled copies on only one of five spans tried.

\cardfield{What goes in, what comes out}
In go the events. The same events come out, each labelled with its family number (1 is the
largest family), or $-1$ with a reason when it fits none. Each family carries a verdict:
whether its members are more alike than stretches of plain background noise from the same
recording would be.

\cardfield{Settings that matter}
\emph{method} picks how families are made: \texttt{hdbscan} (the default) finds tight
crowds and leaves stragglers out; \texttt{kshape} makes a set number of families;
\texttt{ward} puts every event somewhere but runs no noise test.
\emph{hdbscan\_min\_cluster\_size} is the smallest crowd that counts as a family (10); tune
it first. \emph{length\_s} is the stretch cut around each event (600\,s, a placeholder
measured only on synthetic shapes).

\cardfield{Watch out for}
Read each family's noise verdict before its rhythm. In the Intervals table a family marked
\texttt{\textasciitilde noise\_like} or \texttt{\textasciitilde untested} is not a real
family rhythm. Each stretch is rescaled before comparing, so size does not count, only
shape. Events too close to either end of the span for a full stretch are left out.
